\chapter{Prácticas}

Las prácticas siguientes complementan la formación del personal y deben realizarse bajo la supervisión del instructor y con fuentes simuladas cuando así se indique.

\section{Práctica 1: absorción de las radiaciones gamma}

\subsection*{Objetivo}
Reconocer la diferente capacidad de distintos materiales para absorber las radiaciones gamma y determinar cuáles son mejores como blindaje.

\subsection*{Material y equipo}
\begin{itemize}
  \item un detector de radiaciones tipo Geiger--Müller;
  \item un contenedor de gammagrafía con una fuente radiactiva cuya rapidez de exposición en los costados sea aproximadamente de 50 a 100\,mR/h;
  \item tres láminas delgadas de madera (triplay) de $30\times30$\,cm;
  \item tres láminas delgadas de aluminio de $30\times30$\,cm;
  \item tres láminas delgadas de fierro de $30\times30$\,cm.
\end{itemize}

\subsection*{Desarrollo}
Coloque el contenedor en el suelo, con uno de sus costados hacia arriba. Tome y anote en una libreta un mínimo de tres lecturas de rapidez de exposición a contacto. Coloque una lámina de madera unida al costado y tome nuevamente tres lecturas. Repita la secuencia con dos y tres láminas de madera y, por separado, con las láminas de aluminio y de fierro.

Calcule los promedios correspondientes a una, dos y tres láminas de cada material. Compare los resultados obtenidos con una lámina de madera, una de aluminio y una de fierro; formule los comentarios pertinentes y defina cuál parece ser el mejor blindaje contra las radiaciones gamma.

\begin{table}[htbp]
\centering
\caption{Datos de rapidez de exposición con diferentes materiales.}
\begin{tabular}{c c c c}
\hline
 & Madera & Aluminio & Fierro \\
\hline
1 placa & & & \\
2 placas & & & \\
3 placas & & & \\
\hline
\end{tabular}
\end{table}

\section{Práctica 2: manejo y uso del dosímetro de lectura directa}

\subsection*{Objetivo}
Reconocer que los dosímetros de lectura directa son una herramienta de gran utilidad para la seguridad radiológica del radiógrafo.

\subsection*{Material}
Dosímetro de lectura directa por alumno, cargador de dosímetro de lectura directa, contenedor de gammagrafía con una fuente que produzca aproximadamente de 50 a 100\,mR/h a contacto y un cronómetro.

\subsection*{Desarrollo}
Entregue a cada alumno un dosímetro de lectura directa y verifique que su carga sea correcta. Si no marca cero, cárguelo con el cargador mostrado previamente. Colóquelo a un lado o encima del contenedor, en un campo de rapidez de exposición conocida.

Coloque el mayor número posible de dosímetros y controle el tiempo de permanencia de cada uno junto al contenedor. Retire un dosímetro cada 10 minutos y anote su lectura, considerando el factor de calibración. Con los valores obtenidos trace una gráfica que debe coincidir con la gráfica de dosis contra tiempo calculada a partir de la rapidez de exposición medida a un costado del contenedor.

Si la rapidez de exposición es de 100\,mR/h, la dosis acumulada esperada es proporcional al tiempo: 100\,mR al cabo de una hora y 200\,mR al cabo de dos horas. Compare esta referencia con las lecturas de los dosímetros y formule las conclusiones sobre el uso correcto del equipo.

\section{Prácticas 3 y 4: manejo y uso de alarmas sonoras y Geiger--Müller}

\subsection*{Objetivo}
Comprender que estos equipos son herramientas esenciales para trabajar dentro de parámetros de seguridad y para reconocer cuándo se está dentro de un campo de alta rapidez de exposición.

\subsection*{Equipo}
Un detector portátil de radiaciones con alarma sonora, un detector portátil de radiaciones tipo Geiger--Müller y un contenedor de gammagrafía con una fuente cuya rapidez de exposición en los costados sea aproximadamente de 50 a 100\,mR/h.

\subsection*{Alarma sonora}
Encienda el indicador colocando el interruptor en la posición de encendido y seleccione la escala. Aproxímelo casi a contacto del contenedor: debe escucharse claramente el sonido intermitente. Si no se escucha, reemplace la batería verificando que esté correctamente colocada y vuelva a probar el equipo. Si no funciona, notifíquelo de inmediato al instructor.

\subsection*{Detector Geiger--Müller}
Encienda el equipo y coloque la perilla selectora en la posición «batería». Observe que la aguja se desplace a la posición indicada en la carátula. Si no llega a la posición adecuada, cambie las baterías verificando su instalación. Coloque el detector a 50\,cm.

Verifique que la rapidez de exposición corresponda al valor esperado para ese punto y a la actividad de la fuente. Si el equipo no funciona o la lectura difiere en más de 20\% del valor esperado, notifíquelo inmediatamente al instructor para su reemplazo.

\section{Práctica 5: localización de fuentes con un Geiger--Müller}

\subsection*{Objetivo}
Prepararse para encontrar y rescatar una fuente radiactiva, conservando la calma y aplicando correctamente los procedimientos de localización.

\subsection*{Material}
Un contenedor de gammagrafía con una fuente que produzca aproximadamente de 50 a 100\,mR/h en sus costados, un detector portátil de radiaciones tipo Geiger--Müller y un cronómetro.

\subsection*{Desarrollo}
Forme grupos de dos personas. Esconda el contenedor en un patio o área cuya localización sea difícil, sin que los alumnos observen la operación. Permita que el primer grupo inicie la búsqueda con el detector. Los demás participantes observarán la operación para analizar sus aciertos y errores. Registre el tiempo empleado por cada grupo para localizar la fuente.

\section{Práctica 6: aplicación de la ley del cuadrado inverso}

\subsection*{Objetivo}
Aplicar la ley del cuadrado inverso a la rapidez de exposición esperada a partir de una fuente radiactiva colocada en un lugar conocido.

\subsection*{Material}
Un detector portátil de radiación tipo Geiger--Müller y un contenedor de gammagrafía con una fuente cuya rapidez de exposición en los costados sea aproximadamente de 50 a 100\,mR/h.

\begin{figure}[htbp]
\centering
\begin{tikzpicture}[x=0.45cm,y=0.018cm]
  \draw[thick] (0,0) -- (22,0) node[right] {Distancia (cm)};
  \draw[thick] (0,0) -- (0,150) node[above] {Intensidad de radiación};
  \draw[dashed] (0,130) -- (5,130) -- (5,0);
  \draw[dashed] (5,70) -- (10,70) -- (10,0);
  \draw[dashed] (10,30) -- (15,30) -- (15,0);
  \foreach \x in {5,10,15,20} \draw (\x,0) -- (\x,-5) node[below] {\x};
\end{tikzpicture}
\caption{Esquema de la disminución de la intensidad al aumentar la distancia.}
\end{figure}

\section{Práctica 7: atenuación de la radiación gamma}

\subsection*{Objetivo}
Determinar cómo protegerse de las radiaciones utilizando los diferentes blindajes.

\subsection*{Material y desarrollo}
Utilice un detector portátil tipo Geiger--Müller, un contenedor de gammagrafía con una fuente cuya rapidez de exposición sea aproximadamente de 50 a 100\,mR/h y tres láminas delgadas de fierro de $30\times30$\,cm.

Coloque el contenedor en el suelo con un costado hacia arriba y mida la rapidez de exposición a contacto. Anote el valor. Coloque una lámina metálica sobre el costado y repita la medición; repita después con la segunda y la tercera lámina. Con los datos construya una gráfica de rapidez de exposición contra espesor de fierro.

\section{Práctica 8: manejo de un equipo de gammagrafía}

\subsection*{Objetivo}
Al terminar satisfactoriamente esta práctica, el alumno debe estar en posibilidad de armar, utilizar, desarmar y dar mantenimiento adecuado a un equipo completo de gammagrafía.

\subsection*{Material}
Un contenedor de gammagrafía con todos sus accesorios, una fuente simulada, un detector portátil de radiación tipo Geiger--Müller y un juego completo de croquis del equipo y sus componentes.

\subsection*{Desarrollo}
Entregue al alumno el equipo completo de gammagrafía, incluido el detector Geiger--Müller. El alumno efectuará por sí solo los pasos involucrados, desde el armado inicial y el proceso habitual de radiografiado hasta el desmontaje, la limpieza, el mantenimiento y el almacenamiento seguro. Se analizarán las fallas y se encontrará la solución adecuada. Al finalizar, el instructor evaluará los conocimientos y la eficiencia demostrados. Si participan varias personas, la práctica se repetirá para cada alumno.

\section{Esquemas del equipo de gammagrafía}

Los esquemas de las páginas fuente muestran, en orden, el conector de la fuente, el detalle de la chapa de seguridad del contenedor, el mecanismo del manipulador (REEL), el detalle del conector REEL--contenedor y el acople del chicote, y las posiciones de operación radiográfica.

\subsection{Conector de la fuente}
El conector debe quedar correctamente alineado y asegurado antes de iniciar una operación. Verifique las uniones y el estado de sus componentes para evitar una desconexión imprevista.

\subsection{Chapa de seguridad del contenedor}
La chapa mantiene el cierre del contenedor y evita su apertura accidental o no autorizada. Compruebe su funcionamiento durante las inspecciones del equipo.

\subsection{Manipulador REEL y conector}
El manipulador comprende el control, el cable, las mangueras y los acoples que transmiten el movimiento a la fuente. El chicote se conecta al acople del contenedor; todas las uniones deben quedar firmes antes del uso.

\subsection{Operación radiográfica}
En la posición segura, el telemando y la cámara permanecen conectados por las mangueras guía. La fuente se desplaza desde la cámara hasta el punto focal y, terminada la exposición, se retrae nuevamente a la posición segura. El esquema fuente indica una longitud total aproximada de 14.3\,m: 8\,m hacia el telemando y 6\,m de manguera de salida, con un tramo de 2\,m señalado en esta última.

\begin{enumerate}
  \item \textbf{Posición almacenada:} el cable de control está retraído y la fuente permanece dentro del proyector.
  \item \textbf{Fuente en tránsito:} al girar la manivela, el cable desplaza la fuente por la manguera guía.
  \item \textbf{Fuente en el sitio radiográfico:} la fuente alcanza la posición focal para la exposición.
  \item \textbf{Retracción:} después de la exposición, la fuente vuelve al proyector.
\end{enumerate}

\section{Práctica 9: recuperación de fuentes y uso de manipuladores a distancia}

\subsection*{Objetivo}
Confirmar que, al final del curso, el alumno cuenta con los conocimientos para manejar satisfactoriamente una situación de pérdida de control de una fuente radiactiva.

\subsection*{Material}
Un detector tipo Geiger--Müller, un detector portátil con alarma sonora, un contenedor de gammagrafía con accesorios, una fuente simulada, material de acordonamiento (pinzas largas, contenedor de rescate y embudo) y una madrina de plomo.

\subsection*{Desarrollo}
Suponga que una estructura metálica pesada con filos agudos ha caído en la zona de trabajo y ha cortado o roto el chicote o telemando. Localice y blinde la fuente con la madrina de plomo, que debe estar disponible. Acordone y señalice hasta alcanzar lecturas menores de $0.25\,\mathrm{mR/h}$, intensificando la vigilancia del tránsito de público en las zonas aledañas.

Mida la rapidez de exposición en las proximidades de la fuente y anote los valores a contacto del blindaje y a 1 y 10\,m, para estimar el equivalente de dosis necesario para planear el rescate. Sin retirar el blindaje, separe el chicote de la manguera deformada con la herramienta adecuada. Con pinzas de distancia, sostenga el chicote con la fuente hacia abajo e introdúzcala, con un movimiento rápido y seguro, en el contenedor de rescate auxiliándose del embudo. Retire el embudo y coloque la tapa. Estime el equivalente de dosis recibido, compárelo con el dosímetro personal y elabore un informe detallado.

\section{Práctica 10: actividades después de un accidente}

\subsection*{Objetivo}
Poner en práctica las acciones que deben tomarse cuando ocurre un accidente que afecta la seguridad física y radiológica durante el transporte y la operación de equipos radiográficos, analizando sus posibles causas.

\subsection*{Material}
Un detector tipo Geiger--Müller, un detector portátil con alarma sonora, un contenedor de gammagrafía con accesorios, una fuente simulada y material de acordonamiento y señalización.

\subsection*{Desarrollo}
Suponga que ha ocurrido un accidente automovilístico en el que el técnico y el auxiliar no sufrieron daños físicos de consecuencia. A causa del accidente dejaron sin vigilancia el equipo radiográfico y el vehículo durante el tiempo suficiente para que una persona ajena a la empresa se llevara el contenedor con la fuente radiactiva.

En cuanto se advierta la falta del contenedor:
\begin{enumerate}
  \item dé parte de los hechos a las oficinas centrales o al Encargado de Seguridad Radiológica, a la brevedad posible, para recibir instrucciones;
  \item realice una indagación alrededor del lugar donde se supone que fue robada o perdida la fuente;
  \item trasládese al lugar de los hechos con los equipos necesarios para localizar o recuperar la fuente;
  \item realice un muestreo de la zona con grupos de al menos dos personas por cada detector. Efectúe el rastreo cuidadosamente y sin escandalizar a la población;
  \item si la fuente no se localiza en unas cuantas horas, avise a las autoridades locales y solicite su ayuda;
  \item si se localiza fuera del contenedor, acordone la zona, calcule la posible exposición, distribuya las maniobras entre el personal y simule la recuperación para reducir los tiempos antes de intentarla;
  \item informe detalladamente a la Comisión Nacional de Seguridad Nuclear y Salvaguardias, en compañía del encargado de seguridad radiológica;
  \item una vez recuperada la fuente, practique una prueba de fuga al contenedor y tome niveles de radiación para verificar su estado;
  \item redacte un informe pormenorizado y envíelo a la C.N.S.N.S., junto con copias de los informes de prueba de fuga y de verificación de niveles de radiación.
\end{enumerate}
